\subsection{Tokeny, parametry a kontext}
\label{subsec:llm-tokens-parameters-context}

Aby bylo možné text zpracovat neuronovou sítí, musí být nejprve převeden do numerické reprezentace.
Tento proces se skládá ze dvou kroků: tokenizace a embeddingu.

\subsubsection*{Tokenizace}
\label{subsubsec:tokenization}

Tokenizace rozdělí vstupní text na diskrétní jednotky zvané tokeny, přičemž každý token je přiřazen číselnému identifikátoru ze slovníku modelu.
Moderní tokenizace nepracují na úrovni slov ani jednotlivých písmen, ale využívají metodu \textit{subword tokenizace}.
Nejrozšířenějším algoritmem je Byte Pair Encoding (BPE)\cite{sennrich-bpe}, který iterativně slučuje nejčastěji se vyskytující páry bajtů nebo znaků až do dosažení požadované velikosti slovníku.

Výsledkem je, že běžná anglická slova jsou zpravidla zakódována jako jeden token, méně obvyklá slova jako více tokenů.
Ve zdrojovém kódu jsou klíčová slova programovacích jazyků obvykle jednotlivými tokeny, delší identifikátory jsou rozděleny (viz \figref{fig:tokenization-example}).

\begin{figure}[H]
    \centering
    \begin{verbatim}
        Vstup:   std::vector<int> data = {1, 2, 3};
        Tokeny:  ['std', '::', 'vector', '<', 'int', '>', ' data',
                   ' =', ' {', '1', ',', ' 2', ',', ' 3', '};']
    \end{verbatim}
    \caption{Příklad tokenizace výrazu v jazyce C++}
    \label{fig:tokenization-example}
\end{figure}

Z pohledu systému Kelvin má tokenizace přímý praktický dopad.
Podle dokumentace OpenAI platí jako orientační pravidlo, že jeden token odpovídá přibližně čtyřem znakům anglického textu \cite{openai-tokenizer}.
Zdrojový kód je však tokenizován méně efektivně než přirozený jazyk, protože obsahuje speciální znaky, operátory a delší identifikátory, které bývají rozděleny do více tokenů.
Počet tokenů zdrojového kódu je proto zpravidla vyšší než počet jeho slov.
Pro kód s dlouhými identifikátory a hustou syntaxí může být tento poměr ještě výraznější.

\subsubsection*{Embeddingy a poziční kódování}
\label{subsubsec:embedding-positional-encoding}

Každý token ze slovníku je mapován na hustý vektorový prostor pomocí \textit{embedding vrstvy}.
Výsledný vektor $\mathbf{e}_t \in \mathbb{R}^d$ má typicky 768 až 8~192 dimenzí v závislosti na velikosti modelu.
Embedding vrstva je trénovaná společně se zbytkem modelu a kóduje sémantické a syntaktické vlastnosti tokenů: tokeny s podobnou funkcí nebo kontextem mají v tomto prostoru vysokou kosinovou podobnost.

Protože Transformer zpracovává tokeny paralelně bez pořadí, je k embeddingu přidáno \textit{poziční kódování} nesoucí informaci o pozici tokenu v sekvenci.
Původní Transformer\cite{vaswani-attention} použil deterministické sinusové funkce.
Ale moderní modely (LLaMA, Qwen a další) využívají \textit{rotační poziční kódování} (RoPE), které kóduje relativní vzdálenosti tokenů přímo do výpočtu attention a lépe generalizuje na kontexty delší, než model viděl při trénování.

\subsubsection*{Parametry modelu}
\label{subsubsec:model-parameters}

Parametry modelu jsou \textit{trénovatelné váhy}, konkrétně hodnoty matic v attention vrstvách a feed-forward sítích architektury Transformer.
Jejich počet určuje kapacitu modelu, kde větší modely dokáží zachytit složitější vzory a přirozeněji je zobecňovat.
Škálovací zákony pak ukazují, že výkon roste jako mocninná funkce počtu parametrů, objemu trénovaných dat a výpočetního výkonu.
Přičemž pro dosažení optimálního výkonu by počet trénovaných tokenů měl být přibližně dvacetinásobkem počtu parametrů~\cite{hoffmann-chinchilla}.

\subsubsection*{Kontextové okno}
\label{subsubsec:context-window}

Kontextové okno určuje maximální počet tokenů, které model zpracuje v jednom průchodu.
Jedná se tedy o tvrdé omezení, jehož překročení vede ke zkrácení vstupu nebo nutnosti zpracování v oddělených průchodech.
Velikost okna vzrostla s vývojem modelů dramaticky, GPT-3 měl 2~048 tokenů, moderní modely dosahují 128~000 (GPT-4o) až 1~000~000 tokenů (Gemini 2.5 Pro).

Pro analýzu studentských řešení v systému Kelvin jsou relevantní modely s oknem alespoň 16~000 tokenů, které s rezervou pojmou systémový prompt, zadání úlohy i celý odevzdaný zdrojový kód.

\endinput